\section{Methodology}
\label{sec:methodology}


This section presents the architecture and design principles of the
Cognitive Research Memory Engine (CRME). The objective of CRME is to
provide a persistent research memory infrastructure capable of capturing,
organizing, and connecting research activities across multiple sessions
and project stages.


\subsection{CRME Design Principles}

CRME is designed based on four main principles:

\begin{enumerate}

\item \textbf{Persistent Research Memory:}
Research information should remain accessible beyond individual
interaction sessions.

\item \textbf{Semantic Organization:}
Research elements should be represented as connected knowledge
objects rather than isolated documents.

\item \textbf{Provenance Awareness:}
Decisions, artifacts, and results should maintain traceable origins.

\item \textbf{Research Continuity:}
Previous research states should support future activities and reduce
knowledge loss.

\end{enumerate}



\subsection{CRME Architecture}

The overall architecture of CRME is illustrated in
Fig.~\ref{fig:crme_architecture}.


\begin{figure}[t]
\centering
\begin{tikzpicture}[
node distance=4mm,
box/.style={
rectangle,
rounded corners,
draw,
align=center,
font=\scriptsize,
minimum width=24mm,
minimum height=7mm
},
core/.style={
rectangle,
rounded corners,
draw,
very thick,
align=center,
font=\scriptsize\bfseries,
minimum width=34mm,
minimum height=10mm
},
layer/.style={
draw,
dashed,
rounded corners,
inner sep=3mm
},
arrow/.style={
-{Latex[length=1.5mm]},
thick
}
]

\node[box] (user)
{Researcher\\
AI Assistant};

\node[core, below=8mm of user] (kernel)
{CRME Kernel\\
Cognitive Memory Engine};

\node[box, below left=9mm and 8mm of kernel] (memory)
{Memory Engine\\
Research Objects};

\node[box, below=9mm of kernel] (session)
{Session Engine\\
Goals \& Decisions};

\node[box, below right=9mm and 8mm of kernel] (project)
{Project Engine\\
Dependencies\\
Evolution};

\node[box, below=8mm of session] (semantic)
{Semantic Engine\\
Indexing};

\node[box, below=7mm of semantic] (query)
{Query Engine\\
Retrieval};

\node[box, below=7mm of query] (continuity)
{Research Continuity\\
Knowledge Reuse};


\draw[arrow] (user)--(kernel);

\draw[arrow] (kernel)--(memory);
\draw[arrow] (kernel)--(session);
\draw[arrow] (kernel)--(project);

\draw[arrow] (memory)--(semantic);
\draw[arrow] (session)--(semantic);
\draw[arrow] (project)--(semantic);

\draw[arrow] (semantic)--(query);
\draw[arrow] (query)--(continuity);


\node[layer,
fit=(memory)(session)(project),
label={[font=\scriptsize]above:Memory Layer}]
{};


\node[layer,
fit=(semantic)(query),
label={[font=\scriptsize]above:Intelligence Layer}]
{};


\node[layer,
fit=(continuity),
label={[font=\scriptsize]above:Continuity Layer}]
{};


\end{tikzpicture}


\caption{Overall architecture of the Cognitive Research Memory Engine (CRME).}
\label{fig:crme_architecture}
\end{figure}


The framework consists of several modular components:

\begin{itemize}

\item \textbf{Memory Engine:}
Responsible for storing structured research objects, memory packs, and
persistent knowledge states.

\item \textbf{Session Engine:}
Captures interaction sessions, decisions, goals, messages, and generated
artifacts.

\item \textbf{Project Engine:}
Maintains research project structure including milestones,
dependencies, progress states, and project evolution.

\item \textbf{Semantic Engine:}
Provides semantic indexing and relationship discovery between research
objects.

\item \textbf{Query and Retrieval Layer:}
Enables searching and accessing previously stored research knowledge.

\end{itemize}



\subsection{Research Memory Object Model}

Unlike conventional document-based storage systems, CRME represents
research activities as structured memory objects.

Each research object is defined as:

\begin{equation}
O_i=(id_i,type_i,content_i,time_i,source_i,relation_i)
\end{equation}

where $id_i$ represents the object identifier,
$type_i$ defines the object category,
$content_i$ stores the information content,
$time_i$ records temporal information,
$source_i$ represents provenance information,
and $relation_i$ maintains connections with other objects.


The supported research object categories are represented as:

\begin{equation}
Type(O_i)\in\mathcal{T}
\end{equation}


where:

\[
\mathcal{T} =
\{
\textit{Idea},
\textit{Hypothesis},
\textit{Decision},
\textit{Experiment},
\textit{Result},
\textit{Artifact},
\textit{Publication}
\}
\]


This representation enables CRME to preserve the evolution of research
activities from initial ideas to final scientific outputs.



\subsection{Research Lifecycle Modeling}

The research lifecycle model is presented in
Fig.~\ref{fig:research_lifecycle}.


\begin{figure}[t]
\centering
\begin{tikzpicture}[
node distance=7mm,
box/.style={
rectangle,
rounded corners,
draw,
align=center,
font=\scriptsize,
minimum width=22mm,
minimum height=7mm
},
arrow/.style={
-{Latex[length=2mm]},
thick
},
feedback/.style={
-{Latex[length=2mm]},
dashed,
thick
}
]


\node[box](idea){Idea};

\node[box,below=of idea](hypothesis)
{Hypothesis};

\node[box,below=of hypothesis](decision)
{Decision};

\node[box,below=of decision](experiment)
{Experiment};

\node[box,below=of experiment](result)
{Result};

\node[box,below=of result](artifact)
{Artifact};

\node[box,below=of artifact](publication)
{Publication};


\draw[arrow](idea)--(hypothesis);
\draw[arrow](hypothesis)--(decision);
\draw[arrow](decision)--(experiment);
\draw[arrow](experiment)--(result);
\draw[arrow](result)--(artifact);
\draw[arrow](artifact)--(publication);


\draw[feedback]
(publication.east)
.. controls +(25mm,0) and +(25mm,0) ..
(idea.east);


\node[
draw,
dashed,
rounded corners,
fit=(idea)(publication),
inner sep=7mm,
label={[font=\scriptsize]above:
CRME Research Lifecycle Memory}
]{};


\end{tikzpicture}


\caption{Research lifecycle model supported by CRME.}
\label{fig:research_lifecycle}
\end{figure}


Each research activity is stored as a connected memory object, allowing
relationships between hypotheses, decisions, experiments, results, and
artifacts to be maintained.

The lifecycle can be represented as a directed graph:

\begin{equation}
G=(V,E)
\end{equation}

where:

\begin{itemize}

\item $V$ represents research memory objects.

\item $E$ represents semantic or temporal relationships between objects.

\end{itemize}


This graph-based representation enables future retrieval and reasoning
over previous research states.



\subsection{Software Implementation Mapping}

The relationship between conceptual components and implementation
modules is summarized in Table~\ref{tab:implementation_mapping}.


\begin{table}[t]
\caption{Mapping Between CRME Concepts and Software Modules}
\label{tab:implementation_mapping}

\centering
\small

\begin{tabular}{p{0.34\columnwidth}
                p{0.46\columnwidth}}

\toprule
Concept &
Implementation Module
\\
\midrule

Persistent Memory &
memory\_engine.py
\\

Session Continuity &
session\_engine.py
\\

Project Evolution &
project\_engine.py
\\

Semantic Processing &
semantic\_engine.py
\\

Knowledge Query &
query\_engine.py
\\

System Coordination &
crme\_kernel.py
\\

\bottomrule

\end{tabular}

\end{table}



\subsection{Research Continuity Mechanism}

A key feature of CRME is maintaining continuity between separated
research sessions.

When a new session starts, previous memory states can be retrieved and
used as contextual information. Newly generated decisions and artifacts
are integrated into the existing research memory structure.

The continuity process can be described as:

\begin{equation}
M_{t+1}=M_t\cup\Delta M_t
\end{equation}

where $M_t$ represents the current research memory state and
$\Delta M_t$ represents newly generated research information.

This mechanism enables incremental construction of research knowledge
over time.



\subsection{Summary}

The proposed architecture combines memory persistence, semantic
organization, provenance tracking, and project-level evolution in a
single modular framework. The following section evaluates CRME through
experimental scenarios and comparative analysis.

